\section{Conclusion}\label{sec:conclusion}

We have presented BEAMFS v2, a Linux filesystem with an
electromagnetic-resilience recovery contract stated as two
theorems: \cref{thm:soundness-v21} for recovery under bounded
EMR perturbation, and \cref{thm:saturation-v22} for saturation
observability through an auditable
\texttt{UNCORRECTABLE}-flagged journal entry. The contract was
stated after the previous version of this calculus was empirically
falsified by the companion RadFI instrument and was retracted
publicly. The revision retains the v1 mathematical scaffolding
(state space, layout, consistency relation, invariants
\textbf{I1}--\textbf{I5}) and replaces the perturbation family
hypothesis with the EMR threat model of \cref{sec:emr-threat},
covering both stochastic (Family A) and adversarial (Family B)
electromagnetic perturbations bounded by a transverse saturation
boundary at the Reed--Solomon per-subblock correction capacity.

The implementation is 2{,}198 lines of C, 44\% of the 5{,}000-line
audit budget targeted for safety-critical certification regimes,
publicly released under
\texttt{roastercode/beamfs}~\cite{beamfs-impl} with GPG-signed
tagged commits. The empirical section reports byte-level disk
corruption recovered byte-perfect by the INLINE read path with
autonomic in-place repair (\cref{ssec:stage3}), read-path
performance parity with \texttt{ext4} on median and 95th
percentile (\cref{ssec:stage4}), and a stable conformance
fixture canary block over \(\sim\)380{,}000 cache-hit reads with
zero spurious RS correction events.

The present report is, deliberately and explicitly, a
falsifiable scientific statement. The recovery contract is
testable, the implementation is auditable, and any subsequent
counter-example produced by RadFI or an equivalent fault-injection
methodology will lead to a v3 retraction-and-revision cycle
(where v3 designates the next major revision following the
present v2 of this calculus) rather than to a defensive patch. This is the methodological
position we believe a filesystem intended for critical
infrastructure deployment must adopt.

The immediate continuation tasks, in approximate order, are: the
INLINE write path, an active saturation falsification protocol
for \cref{thm:saturation-v22}, the diagnostic of the deferred
RadFI kprobe interaction, an \texttt{xfstests} subset, and
preparation for \texttt{linux-fsdevel} submission. The longer
horizon includes the v3 verified-compilation work and the
extension of the EMR threat model with more targeted Family B
adversarial waveforms.
